\section{Eliminators}
We present some facts/strong intuitions about eliminators which we use in our proofs in \appref{baseline}.

\begin{figure}
   \begin{minipage}[t]{0.7\textwidth}
      \begin{tabularx}{\textwidth}{rL{2.7cm}L{3cm}}
          &\textit{Eliminator}&
          \lowlight{$\Elim\;\Gamma\,A\;K$}
          \\
          $\sigma, \tau ::=$
          &
          $\exVar{e}$
          &
          expression
          \\
          &
          $\elimPhi{\seq{\exPair{p}{\sigma}}}$
          &
          branch node
          \\
          &
          $\emptyset$
          &
          empty eliminator
          \\
          &\textit{Extended Patterns}&
          \\
          $\pstar ::=$
          &
          $\pattConstr{c}{}$
          &
          constructor pattern
          \\
          &
          $\exRecord{}_{\seq{x}}$
          &
          record pattern
          \\
          &
          $\exVar{x}$
          &
          ground variable pattern
          \\
      \end{tabularx}
   \end{minipage}
\caption{Grammar for new eliminators}
\end{figure}

\begin{figure}
   {\small \flushleft \shadebox{$\seq{p},\sigma \elimmapsto \tau$}
   \begin{smathpar}
      \inferrule*
      {
         \strut
      }
      {
         \seq{p},\exVar{e} \elimmapsto \exVar{e}
      }
      %
      \and
      %
      \inferrule*
      {
         \strut
      }
      {
         \seq{p},\emptyset \elimmapsto \emptyset
      }
      %
      \and
      %
      \inferrule*
      {
         \seq{p},\sigma \elimmapsto \tau
      }
      {
         p\cdot\seq{p},\elimPhi{\exPair{p}{\sigma}\cdot\seq{\exPair{p}{\sigma}}'} \elimmapsto \tau
      }
      %
      \and
      %
      \inferrule*
      {
         \exVar{p'}\notin\phiKeys{\elimPhi{\seq{\exPair{p}{\sigma}}}}
      }
      {
         p'\cdot\seq{p}', \elimPhi{\seq{\exPair{p}{\sigma}}} \elimmapsto \emptyset
      }
   \end{smathpar}
   }
   \caption{Eliminator lookup}
\end{figure}

\begin{figure}
   {\small \flushleft \shadebox{$\sigma \disjjoin \exPair{\seq{p}}{\exVar{e}} = \tau$}
   \begin{smathpar}
      \inferrule*
      {
         \strut
      }
      {
         \exVar{e} \disjjoin \exPair{\seq{p}}{\exVar{e'}} = \exVar{e}
      }
      %
      \and
      %
      \inferrule*
      {
         \strut
      }
      {
         \emptyset \disjjoin \exPair{\varepsilon}{\exVar{e}} = \exVar{e}
      }
      %
      \and
      %
      \inferrule*
      {
         \emptyset \disjjoin \exPair{\seq{p}}{\exVar{e}} = \sigma
      }
      {
         \emptyset \disjjoin \exPair{p\cdot\seq{p}}{\exVar{e}} = \elimPhi{\exPair{p}{\sigma}}
      }
      %
      \and
      %
      \inferrule*
      {
         \sigma \disjjoin \exPair{\seq{p}}{\exVar{e}} = \tau
      }
      {
         \elimPhi{\exPair{p}{\sigma} \cdot \seq{\exPair{p}{\sigma}}' } \disjjoin \exPair{p\cdot\seq{p}}{\exVar{e}} = \elimPhi{\exPair{p}{\tau} \cdot \seq{\exPair{p}{\sigma}}'}
      }
      %
      \and
      %
      \inferrule*
      {
         p' \notin \phiKeys{\elimPhi{\seq{\exPair{p}{\sigma}}}}
         \\
         \emptyset \disjjoin \exPair{\seq{p}}{\exVar{e}} = \tau
      }
      {
         \elimPhi{\seq{\exPair{p}{\sigma}}} \disjjoin \exPair{p'\cdot\seq{p}}{\exVar{e}} = \elimPhi{\seq{\exPair{p}{\sigma}} \cdot \exPair{p'}{\tau}}
      }
   \end{smathpar}}
\end{figure}



\begin{theorem}
   \label{thm:join-preserves-join}
   The disjoint join of eliminators is joinability-preserving, that is:
   $$
      \upsigma_1 \disjjoin \upsigma_2 = \uptau' \opAnd \upsigma_2 \disjjoin \upsigma_3 = \uptau'_1 \opAnd \upsigma_1 \disjjoin \upsigma_3 = \uptau'_2
      \newline
      \implies \upsigma_3 \disjjoin \uptau' = \upsigma_2 \disjjoin \uptau'_1 = \upsigma_1 \disjjoin \uptau'_2 = \uptau
   $$
\end{theorem}
\begin{proof}
   Temporarily Omitted.
\end{proof}

An eliminator with no instances of the \textit{branch} constructor, we call it a \textit{singleton} eliminator.
Singletons allow us to define a notion of an eliminators \textit{prefix}.

\begin{definition}
   If $\exists \upsigma'.\; \upsigma \elimmapsto \upsigma' = \uptau$ then we say $\upsigma$ is a prefix of $\uptau$, 
   and write $\upsigma \leq \uptau$. Given two eliminators $\upsigma,\upsigma'$ we write $\uptau = \upsigma \sqcap \upsigma'$
   to mean $\uptau$ is the longest shared prefix of $\upsigma$ and $\upsigma'$
\end{definition}